% Extraído de Textuais/03-Metodologia.tex
% Fonte: git log --follow --name-status do arquivo GitlabRepositorySettingsPage.java (org.eclipse.mylyn, 2026-11-09)
% e 5-Coleta_features/Mylyn/data/method_changes.parquet. Eixo do tempo fora de escala (posições em ordem cronológica).
\begin{figure}[htbp]
    \caption{Nome do método \texttt{create\-Additional\-Controls} do Mylyn e arquivo que o contém ao longo do tempo}
    \label{fig:met-gitlog-nomes}
    \centering
    \begin{tikzpicture}[x=0.92cm,font=\scriptsize,
        ev/.style={circle,draw=black!70,fill=#1,inner sep=2.1pt},
        up/.style={align=center,anchor=south,inner sep=1pt,yshift=3pt},
        dn/.style={align=center,anchor=north,inner sep=1pt,yshift=-3pt},
        snap/.style={black!40,dashed,line width=0.5pt},
        stit/.style={align=center,text=black!60},
        nome/.style={anchor=south west,align=left,fill=white,inner sep=1.5pt}]
        % releases (snapshots)
        \foreach \x/\r/\d in {0.0/41/jun. 2023, 6.05/42/set. 2023, 8.6/43/nov. 2023, 12.2/44/mar. 2024} {
            \draw[snap] (\x,-0.95) -- (\x,5.75);
            \node[stit,above] at (\x,5.75) {\textit{release} \r\\\d};
        }
        \node[stit] at (13.0,6.05) {\ldots};
        % faixa do nome antigo
        \draw[line width=2.6pt,black!25] (0.6,3.0) -- (3.0,3.0);
        \node[nome] at (0.24,4.33)
            {\textbf{Nome antigo} \texttt{}};
        % faixa do nome novo
        \draw[-{Stealth[length=6pt]},line width=2.6pt,black!25] (3.6,0) -- (16.3,0);
        \node[nome] at (4.0,0.95)
            {\textbf{Nome novo} \texttt{}};
        % mudança de pacote
        \draw[-{Stealth[length=6pt]},line width=1.1pt,blue!60!black] (3.0,3.0) to[out=-70,in=160] (3.6,0);
        \node[align=center,font=\scriptsize\itshape,text=blue!60!black] at (1.75,1.4) {mudança\\de nome};
        % commits no nome antigo
        \node[ev=red!45] at (0.9,3.0) {};   \node[up] at (0.9,3.0) {\texttt{b53db5e05c}\\2023-06-07\\cria o arquivo};
        \node[ev=red!45] at (2.0,3.0) {};   \node[dn] at (2.0,3.0) {\texttt{011737ce02}\\2023-07-17};
        \node[ev=blue!35] at (3.0,3.0) {};  \node[up] at (3.0,3.0) {\texttt{9efc5a0011}\\2023-07-24};
        % commits no nome novo
        \foreach \x/\c/\d/\cor/\pos in {4.3/d0a2f4d4f0/2023-07-28/ForestGreen!55/up, 5.3/6574dcfa14/2023-09-03/yellow!60/dn,
                                       6.9/f3d06000be/2023-10-03/ForestGreen!55/up, 7.9/1a1e9bb1fb/2023-11-05/yellow!60/dn,
                                       9.4/b9e87cf498/2024-02-02/yellow!60/up, 10.5/180f862c04/2024-02-05/yellow!60/dn,
                                       11.6/bf80278d6e/2024-02-17/yellow!60/up, 13.4/b9a87ac56d/2025-07-20/yellow!60/dn,
                                       14.6/eef8dc9550/2025-12-27/yellow!60/up} {
            \node[ev=\cor] at (\x,0) {};
            \node[\pos] at (\x,0) {\texttt{\c}\\\d};
        }
        \node[stit,anchor=north] at (15.7,-0.15) {versão final\\2026};
        % legenda
        \begin{scope}[shift={(0.0,-1.55)}]
            \node[ev=yellow!60] at (0,0) {};        \node[anchor=west] at (0.15,0) {nas duas listas};
            \node[ev=ForestGreen!55] at (3.0,0) {}; \node[anchor=west] at (3.15,0) {só na coleta};
            \node[ev=red!45] at (0,-0.45) {};       \node[anchor=west] at (0.15,-0.45) {só no \texttt{git log -L} (a coleta contou \texttt{011737ce02} para o nome antigo)};
            \node[ev=blue!35] at (0,-0.9) {};       \node[anchor=west] at (0.15,-0.9) {toca o arquivo, mas em nenhuma das listas};
        \end{scope}
    \end{tikzpicture}
    \fonte{Elaborado pelo autor (2026).}
\end{figure}%
